% Formulario para que un chofer o cobrador ejerza sus derechos (acceso,
% rectificación, cancelación, oposición, portabilidad) o retire su
% consentimiento. Lo recibe la cooperativa, que es la responsable; el
% sistema tiene con qué contestar cada pedido (ver la tabla del final).
\documentclass[10.5pt,a4paper]{article}
% Preámbulo común de los documentos legales (legal/*.tex).
% Se compila cada uno con:   cd legal && latexmk -pdf <archivo>.tex
% La paleta es la del producto y la del manual: si cambia una, cambian las dos.
\usepackage[T1]{fontenc}
\usepackage[utf8]{inputenc}
\usepackage[spanish,es-noshorthands,es-tabla]{babel}
\usepackage[scaled=0.92]{helvet}
\renewcommand{\familydefault}{\sfdefault}
\usepackage{amssymb}
% Los que usan los textos y pdflatex no trae
\DeclareUnicodeCharacter{2713}{\checkmark}
\DeclareUnicodeCharacter{22EF}{\ensuremath{\cdots}}
\DeclareUnicodeCharacter{2192}{\ensuremath{\rightarrow}}
\usepackage[a4paper,top=2cm,bottom=2.2cm,left=2.2cm,right=2.2cm,headheight=14pt]{geometry}
\usepackage[table]{xcolor}
\usepackage{microtype}
\usepackage{enumitem}
\usepackage{booktabs}
\usepackage{tabularx}
\usepackage{array}
\usepackage{longtable}
\usepackage{fancyhdr}
\usepackage{titlesec}
\usepackage{lastpage}
\usepackage[most]{tcolorbox}
\usepackage[hidelinks]{hyperref}

\definecolor{tinta}{HTML}{0B2239}
\definecolor{tinta2}{HTML}{2A4864}
\definecolor{tenue}{HTML}{496A86}
\definecolor{marca}{HTML}{1668C4}
\definecolor{ambar}{HTML}{A67300}
\definecolor{rojo}{HTML}{C2001D}
\definecolor{fondo}{HTML}{F4F7FB}
\definecolor{linea}{HTML}{C3D4E4}
\definecolor{ambarfondo}{HTML}{FDF6E4}
\definecolor{ambarborde}{HTML}{E4C88A}
\definecolor{azulfondo}{HTML}{E7F0FB}
\definecolor{azulborde}{HTML}{C9DDF4}

\color{tinta}
\setlength{\parindent}{0pt}
\setlength{\parskip}{5pt plus 2pt}
\linespread{1.08}
\setlist{itemsep=2pt,topsep=3pt}

\titleformat{\section}{\color{tinta}\large\bfseries}{}{0em}{}
\titleformat{\subsection}{\color{tinta2}\normalsize\bfseries}{}{0em}{}
\titlespacing*{\section}{0pt}{12pt}{4pt}
\titlespacing*{\subsection}{0pt}{8pt}{2pt}

\renewcommand{\arraystretch}{1.25}
\arrayrulecolor{linea}

% Lo que hay que completar a mano: una línea para escribir
\newcommand{\llenar}[1][6cm]{\rule[-2pt]{#1}{0.5pt}}
% Lo que tiene que decidir o completar el abogado o las partes
\newcommand{\completar}[1]{{\color{ambar}[#1]}}
% Una casilla para marcar
\newcommand{\casilla}{\fbox{\rule{0pt}{7pt}\rule{7pt}{0pt}}}

\newtcolorbox{nota}[1][]{colback=azulfondo,colframe=azulborde,boxrule=0.6pt,arc=3pt,
  left=8pt,right=8pt,top=5pt,bottom=5pt,fontupper=\small\color{tinta2},#1}
\newtcolorbox{borrador}{colback=ambarfondo,colframe=ambarborde,boxrule=0.6pt,arc=3pt,
  left=8pt,right=8pt,top=5pt,bottom=5pt,fontupper=\small}

% Pie: qué documento es, su versión y la página
\newcommand{\documento}[2]{%
  \pagestyle{fancy}\fancyhf{}%
  \renewcommand{\headrulewidth}{0pt}%
  \fancyfoot[L]{\scriptsize\color{tenue}#1 · #2}%
  \fancyfoot[R]{\scriptsize\color{tenue}Página \thepage\ de \pageref*{LastPage}}%
}

% Firma: una línea y lo que va debajo
\newcommand{\firma}[1]{\begin{minipage}[t]{0.44\textwidth}\vspace{1.6cm}\rule{\linewidth}{0.5pt}\\{\small\color{tinta2}#1}\end{minipage}}

\begin{document}
\documento{Solicitud de derechos sobre datos personales}{borrador del 11/11/2026}

{\LARGE\bfseries Solicitud sobre mis datos personales}\par
{\small\color{tenue}Ley N.º 29733 · Aplicación «Control de ruta» · se entrega a la cooperativa (a su central)}

\begin{borrador}
\textbf{Borrador.} Que lo revise el abogado junto con el consentimiento. La
cooperativa tiene que contestar dentro de los plazos de ley
\completar{confirmar: acceso 20 días hábiles; rectificación, cancelación y
oposición 10 días hábiles}.
\end{borrador}

\section{Quién lo pide}
Nombre completo: \llenar[11cm]\par\vspace{5pt}
DNI: \llenar[4cm]\quad Usuario en la app: \llenar[4.5cm]\quad Rol: \llenar[3cm]\par\vspace{5pt}
Teléfono o correo para la respuesta: \llenar[9cm]

\section{Qué pide}
Marque lo que pide. Puede marcar más de una.

\begin{tabularx}{\linewidth}{@{}p{14pt}X@{}}
\casilla & \textbf{Acceso}: saber qué datos míos tiene el sistema.\\
\casilla & \textbf{Portabilidad}: una copia de todos mis datos en un archivo.\\
\casilla & \textbf{Rectificación}: corregir un dato que está mal (diga cuál y cómo debe quedar).\\
\casilla & \textbf{Cancelación}: que se borren mis datos que no sean registro de trabajo de la cooperativa.\\
\casilla & \textbf{Oposición}: que no se usen mis datos para un fin concreto (diga cuál).\\
\casilla & \textbf{Revisión de una marca automática} (por ejemplo «GPS sospechoso» o «entrada tardía»): que una persona la revise y me la explique.\\
\casilla & \textbf{Retiro del consentimiento} que firmé (versión \llenar[2.6cm]). Entiendo que la ubicación durante el turno es la función central de la aplicación y que la cooperativa me dirá cómo se registra mi trabajo sin ella.\\
\end{tabularx}

\section{Detalle}
\vspace{2pt}
\llenar[\linewidth]\par\vspace{8pt}\llenar[\linewidth]\par\vspace{8pt}\llenar[\linewidth]\par\vspace{8pt}\llenar[\linewidth]

\vspace{6pt}
Lugar y fecha: \llenar[8cm]

\noindent\firma{Firma de quien lo pide}\hfill\firma{Recibido por la cooperativa (nombre, cargo, fecha y firma)}

\vspace{14pt}
\begin{nota}
\textbf{Para la cooperativa: cómo se contesta con el sistema.}
\begin{itemize}[leftmargin=12pt,itemsep=0pt]
\item \textbf{Acceso y portabilidad}: Despacho → Gestión → Personas → «⋯» →
  Editar → \emph{Bajar todos sus datos (ZIP)}. Trae un LÉEME que explica cada
  archivo. Queda en la Actividad quién lo bajó.
\item \textbf{Rectificación}: Personas → «⋯» → Editar (nombre y alias). Las
  horas se corrigen en Turnos, con motivo.
\item \textbf{Cancelación}: Personas → «⋯» → Baja. Borra la cuenta, sus combis
  habilitadas y sus mensajes privados. Los turnos, vueltas y la Actividad son
  registro de trabajo y se quedan hasta su plazo \completar{confirmar con el abogado}.
\item \textbf{Retiro del consentimiento}: Personas → Consentimiento → «Lo
  retiró», con cómo lo pidió; y en la hoja de registro.
\item \textbf{Marcas automáticas}: se miran en Gestión → Dónde se traba y en
  Números → Señal y presencia, con su detalle y hora.
\end{itemize}
\end{nota}
\end{document}
